\documentclass[12pt,a4paper]{article}

\usepackage[utf8]{inputenc}
\usepackage[T1]{fontenc}
\usepackage{amsmath}
\usepackage{graphicx}
\usepackage{booktabs}
\usepackage{geometry}
\usepackage{hyperref}
\geometry{margin=2.5cm}

\title{Understanding Cat Behavior:\\A New Owner's Guide to Earning Trust}
\author{Ittipol}
\date{\today}

\begin{document}

\maketitle

\begin{abstract}
Bringing a cat into a new home is a rewarding experience, but it can also be confusing
for a first-time owner. Cats communicate through subtle body language, and many behaviors
that seem strange or even hostile are simply a cat's way of expressing uncertainty.
This article explains the most common feline behaviors a new owner will encounter,
decodes their meaning, and offers a practical, step-by-step approach to building trust
with a newly adopted cat. With patience, consistency, and an understanding of feline
psychology, most cats transition from fearful newcomers to affectionate companions
within a few weeks.
\end{abstract}

\tableofcontents
\newpage

\section{Introduction}
Cats are one of the most popular companion animals in the world, yet they are also one
of the most misunderstood. Unlike dogs, cats are solitary hunters by descent, and their
social bonds are built on choice rather than hierarchy. A cat that hides, stares, or
refuses to be touched is not being ``mean''---it is behaving exactly as evolution
designed it to respond to uncertainty. Understanding this perspective is the first step
toward earning a cat's trust.

This guide is organized in three parts: reading cat body language (Section~2),
interpreting common behaviors (Section~3), and a practical trust-building program
for new owners (Section~4).

\section{Reading Cat Body Language}
A cat's emotional state is written in its posture, tail, ears, and eyes. Table~\ref{tab:body}
summarizes the most important signals.

\begin{table}[h]
\centering
\caption{Common cat body-language signals and their meaning.}
\label{tab:body}
\begin{tabular}{@{}lll@{}}
\toprule
\textbf{Signal} & \textbf{Appearance} & \textbf{Meaning} \\
\midrule
Tail held high & Tail vertical, slight curve at tip & Confident, friendly \\
Purring & Rhythmic low vibration & Content (or self-soothing) \\
Slow blink & Eyes half-closed, relaxed & Trust, non-threat \\
Ears flattened & Ears turned sideways or back & Fear or irritation \\
Dilated pupils & Large black eyes in bright light & Fear, excitement, or play \\
Exposed belly & Cat rolls onto its back & Deep relaxation (not always an invite!) \\
Hiding & Cat retreats under furniture & Overwhelmed; needs space \\
\bottomrule
\end{tabular}
\end{table}

One of the most useful signals for a new owner is the \emph{slow blink}. In feline
communication, sustained staring is a threat, while slow, deliberate blinking signals
safety. Returning a slow blink is one of the simplest ways to tell a nervous cat that
you are not a danger.

\section{Common Behaviors Explained}
\subsection{Hiding}
Hiding is normal for the first several days in a new home. A cat's instinct tells it
that unknown territory may contain predators. Forcing a cat out of hiding reverses
trust-building; letting it emerge on its own schedule accelerates it.

\subsection{Kneading (``Making Biscuits'')}
Rhythmic pushing of the paws against a soft surface is a kitten behavior carried into
adulthood. It indicates contentment and comfort---often performed on a trusted person's
lap.

\subsection{Head Bunting and Rubbing}
Cats have scent glands on their cheeks and forehead. Rubbing against you marks you as
part of their social group and is one of the clearest signs of acceptance.

\subsection{Scratching Furniture}
Scratching is not destructive spite; it maintains claw health and marks territory
visually and by scent. Providing a scratching post near the cat's favorite spot
redirects the behavior without conflict.

\section{Building Trust: A Step-by-Step Program}
Trust develops in stages. Attempting a later stage before the cat is ready for an
earlier one sets progress back. The suggested timeline in Table~\ref{tab:trust}
fits most newly adopted cats.

\begin{table}[h]
\centering
\caption{A staged approach to earning a cat's trust.}
\label{tab:trust}
\begin{tabular}{@{}clp{8.5cm}@{}}
\toprule
\textbf{Days} & \textbf{Stage} & \textbf{What to do} \\
\midrule
1--3 & Decompress & Provide a quiet room with food, water, litter, and a hiding place. Visit briefly; do not force contact. \\
3--7 & Associate & Sit in the room, speak softly, toss treats near the cat. Let it approach you. \\
1--2 weeks & First contact & Offer a finger to sniff; try gentle chin or cheek scratches if invited. Use slow blinks. \\
2--4 weeks & Play & Interactive toys (wand teasers) let the cat build confidence at a distance. \\
1--3 months & Bond & Predictable feeding times, gentle grooming, and lap invitations on the cat's terms. \\
\bottomrule
\end{tabular}
\end{table}

\subsection{A Simple Model of Trust Growth}
The progress of trust-building is not linear---early interactions matter enormously,
while later ones reinforce an established bond. A toy model captures this idea with a
saturating curve:
\begin{equation}
T(t) = T_{\max}\left(1 - e^{-kt}\right),
\label{eq:trust}
\end{equation}
where $T(t)$ is the level of trust after $t$ days, $T_{\max}$ is the cat's maximum
potential trust, and $k$ depends on the consistency of positive interactions such as
feeding routines, play sessions, and calm body language. A new owner cannot force
$T$ to rise faster than the cat's nature allows; they can only raise $k$ through
patience and predictability.

\section{Conclusion}
A cat's behavior is a readable language once the owner knows what to look for. Signals
such as the slow blink, the high tail, and kneading tell an owner when the relationship
is deepening, while hiding and flattened ears honestly report that the cat needs more
time. Trust is earned in stages---decompression, association, first contact, play, and
bond---and cannot be rushed. An owner who respects the cat's timeline is usually
rewarded with years of quiet, genuine companionship.

\begin{thebibliography}{9}
\bibitem{bradshaw} J. Bradshaw, \emph{Cat Sense: How the New Feline Science Can Make
You a Better Friend to Your Pet}. Basic Books, 2013.
\bibitem{johnson} P. Johnson-Bennett, \emph{Think Like a Cat: How to Raise a
Well-Adjusted Cat---Not a Sour Puss}. Penguin Books, 2011.
\bibitem{icatcare} International Cat Care, ``Bringing a new cat home.''
\url{https://icatcare.org}
\bibitem{aspca} ASPCA, ``Cat behavior.'' \url{https://www.aspca.org}
\end{thebibliography}

\end{document}
